\section{Problem and contribution}
\label{sec:introduction}

The cost of a coding agent includes repeated repository discovery, specification
interpretation, planning, repair, validation, and recovery. Persistent state can
avoid some of this work only when the state remains applicable to the current
program. A shorter prompt can omit necessary dependencies; a cached pass can
refer to obsolete source. Repository-agent interfaces and test-validation
benchmarks make this boundary between proposals and checks central to useful
repair~\citep{baseline08,baseline11}.

Our proposed method has three connected contributions. First, source and intent
produce identity-bound semantic facts and explicit unresolved obligations.
Second, those obligations guide deterministic analysis, checked reasoning and
bounded synthesis before an authorized neural residual. Third, validation,
evidence reuse and accepted-state publication preserve the scope and
applicability of each observation. Accepted changes invalidate affected facts
and plans, providing the next iteration with persistent, current evidence.

The resulting hypothesis concerns the entire workflow: repeated neural
reasoning and validation may be displaced by persistent, checked work.
Context compression is one possible mechanism. Fewer model calls and less
replanning are distinct mechanisms, while test reuse primarily reduces
validation computation. Each requires its own correctness and cost comparison.
No benefit follows merely from having a symbolic component in the source tree.

We retain an earlier additive-context cold-repair study as one bounded
experiment, then evaluate missing mechanisms incrementally. Architecture,
local qualification, development pilots and confirmatory effectiveness have
different evidential roles. A failed or inconclusive experiment remains a
result; an unrun comparison remains an open research obligation. In particular,
the original 32 cold cells cannot establish the value of persistent reasoning
or whole-system token savings.
